\part{Architecture logicielle, diagnostic et travaux pratiques}
\section{Chaîne d'encodage et chaîne de décodage}
\subsection{Encoder une image développée}
Un export doit définir son contrat avant de choisir un
codec : dimensions, orientation, espace couleur, fonction
de transfert, profondeur, alpha et métadonnées. Une chaîne
conceptuelle correcte est : terminer les traitements dans
leur domaine approprié, effectuer le rendu de sortie,
convertir les couleurs, appliquer l'encodage de transfert,
quantifier, compresser puis écrire les structures et les
métadonnées cohérentes.

Le codec peut effectuer lui-même certaines de ces opérations.
Il faut connaître son API pour ne pas les appliquer deux
fois. Une conversion RGB vers YCbCr peut être gérée par
l'encodeur JPEG ; une bibliothèque TIFF peut attendre
des octets déjà organisés comme les tags le décrivent.

L'écriture devrait utiliser un fichier temporaire dans
un emplacement approprié, vérifier les erreurs puis
publier le résultat complet. Une taille non nulle ne
prouve pas que l'écriture est terminée. En cas de conflit
de nom, une stratégie d'incrémentation doit tester le
nom final complet et préserver les fichiers existants.
L'atomicité exacte dépend du système de fichiers et de
l'opération de publication choisie.

\subsection{Décoder une image rendue}
Le lecteur vérifie la structure, choisit le codec,
reconstruit les composantes, interprète leur type,
leur couleur et leur alpha, applique l'orientation
si demandée et produit un tampon interne. La conversion
vers l'écran vient ensuite. Il faut savoir si la
bibliothèque réalise déjà orientation, gestion ICC,
conversion de profondeur ou prémultiplication.

Un contrat comme « RGBA float32 linéaire sRGB, alpha droit,
orientation appliquée » est beaucoup plus utile que
« une image OpenCV ». Les erreurs inter-modules viennent
souvent d'informations restées implicites.

\section{Mathématiques de traitement : une carte pour se repérer}
\subsection{Convolution et séparation fréquentielle}
Un filtre discret linéaire peut s'écrire
\begin{equation}
J[x,y]=\sum_{u,v}K[u,v]I[x-u,y-v].
\end{equation}
Un noyau de moyenne réduit des variations rapides et
du bruit, mais aussi des détails. Un noyau passe-haut
réagit aux contours et au bruit. Les conditions aux
bords, répétition, symétrie ou remplissage, font partie
de l'algorithme et peuvent créer des différences de test.

La transformée de Fourier discrète 2D est
\begin{equation}
F[k,\ell]=\sum_{x=0}^{W-1}\sum_{y=0}^{H-1}
I[x,y]e^{-2\pi i(kx/W+\ell y/H)}.
\end{equation}
Elle décompose des motifs en fréquences globales ;
les ondelettes ajoutent une localisation et des
échelles. FFT désigne un algorithme rapide pour
calculer une transformée, pas un filtre en soi.

\subsection{Accentuation et débruitage}
Une accentuation classique utilise
$J=I+a(I-G_\sigma*I)$ : elle ajoute une différence
entre image et version floutée. Le paramètre $a$
augmente aussi le bruit et peut créer des halos.
L'ordre débruitage/netteté change donc le résultat.

Le débruitage est un problème d'estimation : retrouver
une image plausible à partir de mesures bruitées.
Un objectif simplifié est
\begin{equation}
\widehat I=\arg\min_J\left[
\frac{1}{2\sigma^2}\|Y-J\|_2^2+\lambda\mathcal R(J)\right].
\end{equation}
Le premier terme attache la solution aux données,
le second impose une régularité ou un a priori.
Les méthodes non locales, bayésiennes ou apprises
diffèrent notamment par cet a priori et leur modèle
de bruit. Un bruit de photons dépend du signal ;
un modèle gaussien constant n'est pas toujours suffisant.
Les cours 05, 08 et 10 développent les outils liés.

\subsection{Géométrie et rééchantillonnage}
Une transformation affine s'écrit en coordonnées
homogènes :
\begin{equation}
\begin{pmatrix}x'\\y'\\1\end{pmatrix}=
\begin{pmatrix}a&b&t_x\\c&d&t_y\\0&0&1\end{pmatrix}
\begin{pmatrix}x\\y\\1\end{pmatrix}.
\end{equation}
Une homographie remplace aussi la dernière ligne,
puis nécessite une division par la troisième
coordonnée. Pour produire une grille de sortie
sans trous, on utilise souvent la transformation
inverse : pour chaque pixel cible, chercher sa
position source et interpoler.

Un recadrage entier peut simplement changer les
bornes. Une rotation arbitraire, une correction
de perspective ou de distorsion impose généralement
un rééchantillonnage. Les conventions de centre
de pixel et le filtre doivent être identiques
pour comparer deux implémentations au pixel près.

\section{Aperçus, proxies, caches et performances}
\subsection{Pourquoi une image peut apparaître en deux temps}
Un aperçu intégré est rapide à afficher, mais représente
un rendu déjà calculé, parfois dans une résolution
et un profil différents. Le moteur peut ensuite
remplacer cet aperçu par son développement. Ce
changement ne signifie pas nécessairement qu'un
filtre utilisateur a été appliqué.

Un proxy est une représentation réduite choisie
pour l'interaction. Il peut accélérer le zoom
global et les réglages, mais ne suffit pas à
juger du bruit ou de la netteté à l'échelle native.
Un zoom 100\,\% du proxy n'est pas le 100\,\%
du RAW pleine résolution.

\subsection{Ce qu'une clé de cache doit distinguer}
Une clé robuste inclut l'identité de la source,
le niveau de résolution, l'orientation, les
paramètres de décodage, la calibration ou son
identité, la recette de rendu et la version
de l'étape concernée. Sinon, changer de profil
peut réutiliser une image développée avec l'ancien.

Un cache borné LRU élimine les entrées les moins
récemment utilisées selon un budget explicite.
Le budget doit être en octets ou en coût estimé,
pas seulement en nombre d'images de tailles
potentiellement différentes. Précharger des
vignettes ne justifie pas de garder indéfiniment
tous les RAW développés.

\subsection{Mesurer l'ouverture}
Chronométrer séparément : recherche des profils,
lecture, décompression, dématriçage, transformations,
copie mémoire, transfert GPU et affichage. Comparer
le premier passage, avec caches froids, aux suivants,
avec caches chauds. Une mesure globale de 30 secondes
ne permet pas de désigner le dématriçage ou le disque
comme responsable sans autre observation.

Les en-têtes et métadonnées peuvent être lus sans
développement complet. Une vignette peut provenir
du fichier ou d'un cache. L'architecture doit
prévoir annulation et priorité des demandes :
une image quittant la vue ne doit pas monopoliser
le calcul au détriment de celle que l'on ouvre.

\section{Sécurité, robustesse et archivage}
\subsection{Un fichier image est une entrée non fiable}
Valider tailles, offsets, profondeur, nombre de
canaux, ratio de décompression et profondeur
d'imbrication. Limiter mémoire et durée de calcul.
Une bombe de décompression est un fichier compact
qui demande énormément de mémoire après expansion.
Un graphe de répertoires peut contenir une boucle.
Un profil embarqué peut lui aussi être malformé.

Les CRC n'empêchent pas un attaquant de fabriquer
un fichier cohérent et malveillant. Les bibliothèques
doivent être maintenues, testées sur corpus variés
et éventuellement isolées selon le contexte. Le
fuzzing explore des entrées modifiées pour trouver
des plantages ou comportements indéfinis ; il
complète les tests de conformité, sans les remplacer.

\subsection{Vie privée et provenance}
Un export peut révéler lieu, date, numéro de série,
nom de propriétaire et parfois davantage de contexte
qu'attendu. Une option de suppression des métadonnées
doit distinguer données privées et informations
nécessaires à l'affichage, telles qu'un profil ICC.
Supprimer aveuglément tous les tags peut modifier
l'interprétation de l'image.

\subsection{Archivage durable}
Conserver original, sauvegardes vérifiées, hashes,
recette, profils autorisés et versions significatives.
Un rendu TIFF peut préserver une apparence stabilisée,
mais ne remplace pas toutes les possibilités du RAW.
Un DNG peut améliorer l'interopérabilité sans garantir
que chaque métadonnée constructeur aura été conservée
ou interprétée de manière identique. Tester une
restauration est plus probant que constater qu'une
copie existe.

\section{Diagnostic : quels nombres comparer ?}
\subsection{Trois niveaux de test}
\begin{enumerate}
\item \textbf{Binaire :} signature, structure, offsets,
CRC éventuels, hash et longueur. Cela vérifie le fichier.
\item \textbf{Échantillons :} dimensions, type,
orientation, tableau décodé et round-trip. Cela vérifie
les données reconstruites.
\item \textbf{Rendu :} profil, fonction de transfert,
géométrie, luminance, couleurs et régions. Cela vérifie
l'interprétation et les traitements.
\end{enumerate}
Un round-trip sans perte doit conserver exactement
les échantillons du domaine déclaré. Un rendu RAW
avec dématriçage n'est pas un round-trip réversible
vers la mosaïque d'origine.

\subsection{MAE, RMSE et PSNR}
Pour $N$ échantillons comparables $a_i,b_i$,
\begin{align}
\mathrm{MAE}&=\frac1N\sum_i|a_i-b_i|,\\
\mathrm{RMSE}&=\sqrt{\frac1N\sum_i(a_i-b_i)^2},\\
\mathrm{PSNR}&=20\log_{10}\frac{L_{\max}}{\mathrm{RMSE}}.
\end{align}
La plage $L_{\max}$ et le domaine, codes sRGB ou
valeurs linéaires par exemple, doivent être précisés.
Une erreur nulle donne un PSNR infini dans cette
définition ; il faut le représenter explicitement
dans les rapports plutôt que provoquer une division
par zéro.

Un décalage d'un pixel sur un contour peut produire
une forte MAE sans grande différence de couleur
locale. Une moyenne globale peut masquer une
dégradation importante sur la peau ou les ombres.
Mesurer des régions fixes et visualiser les erreurs.

\subsection{Luminance et erreur colorimétrique}
La luminance relative $Y$ se calcule à partir des
composantes linéaires avec les poids de l'espace,
pas directement avec les codes non linéaires.
Pour des moyennes strictement positives,
\begin{equation}
\Delta e_Y=\log_2\frac{\overline Y_{\text{référence}}}
{\overline Y_{\text{test}}}
\end{equation}
décrit un rapport en EV. Cela ne prouve pas qu'un
curseur d'exposition suffira à corriger le rendu.

La distance Euclidienne CIELAB, notée $\Delta E_{76}$,
n'est pas uniformément perceptuelle. CIEDE2000
pondère et couple différemment les écarts ; sa
mise en œuvre doit être testée avec des paires
de référence, comme celles de Sharma et collègues~\cite{ciedeimages}.
Un $\Delta E$ exige une conversion colorimétrique
cohérente et ne constitue pas seul une évaluation
du détail ou du bruit.

\subsection{Clipping et fractions sombres}
« Au moins un canal proche de zéro » n'est pas
« pixel noir ». Une couleur rouge pure possède
deux canaux nuls sans être noire. Il faut distinguer
fraction de canaux, fraction de pixels avec un
canal écrêté, fraction tous canaux écrêtés et
seuil de luminance. Pour une sortie 16 bits,
indiquer explicitement un seuil comme $1/65535$.

Le bruit estimé par MAD est la médiane des écarts
absolus à une médiane. Sur un coefficient gaussien
centré, $\sigma\approx\mathrm{MAD}/0{,}67448975$.
Il faut choisir une zone ou une transformation
qui sépare suffisamment le bruit de la texture.
Une faible dispersion après flou peut seulement
signifier que l'on a détruit du détail.

\section{Travaux pratiques reproductibles}
\subsection{Lancer le laboratoire autonome}
Depuis la racine du dépôt, avec Python 3 :
\begin{verbatim}
python docs/scholar/references/images-exemples.py
\end{verbatim}
Sans option, il affiche des calculs et lance ses
autotests, sans créer de fichier. Il n'utilise
que la bibliothèque standard Python. Pour fabriquer
trois images synthétiques dans un nouveau dossier :
\begin{verbatim}
python docs/scholar/references/images-exemples.py \
  --output-dir build-scholar-images-tp
\end{verbatim}
La continuation par antislash ci-dessus est celle
d'un shell POSIX ; sous PowerShell, écrire la
commande sur une ligne ou utiliser un accent grave.
Le dossier existant est refusé. Les photographies
du corpus et les profils Adobe ne sont pas modifiés.

Les trois fichiers montrent les mêmes deux pixels
rouge et vert dans PPM, PNG et TIFF. Le programme
compare les pixels de son PNG avec ceux d'origine,
teste les cinq filtres et vérifie les offsets TIFF.
Son décodeur PNG est volontairement limité à RGB
8 bits non entrelacé et refuse les autres variantes.
Il ne faut pas le brancher sur la production.

\subsection{Observer les octets}
\begin{verbatim}
python docs/scholar/references/images-exemples.py \
  --inspect build-scholar-images-tp/deux-pixels.tiff
\end{verbatim}
Cette option lit seulement les 64 premiers octets
et la taille. Comparer la signature au tableau
TIFF de la partie IV. Ouvrir ensuite les trois
fichiers dans un lecteur compatible et vérifier
qu'ils montrent bien deux pixels : rouge puis vert.
Une application peut agrandir ces pixels par
interpolation ; pour observer leurs codes, préférer
une lecture numérique ou un zoom sans lissage.

\subsection{Lire les métadonnées d'une photographie réelle}
Avec ExifTool installé, en lecture seule :
\begin{verbatim}
exiftool -G1 -a -s photo.CR3
exiftool -G1 -a -s photo.tiff
\end{verbatim}
Adapter l'extension à Sony ou Nikon. Relever les
dimensions, modèle, ISO, orientation, profondeur,
profil, compression, gains et informations de
développement disponibles. Les doublons sont
volontairement affichés avec leurs groupes.
Le cours ne recommande aucune commande d'écriture
de métadonnées sur les originaux.

\subsection{Éprouver la réversibilité}
Le laboratoire compacte $[\texttt{ABC},\texttt{123}]$
sur 12 bits et obtient \texttt{AB C1 23}. Deux
échantillons font 24 bits, donc trois octets.
Modifier un seul bit et observer quelle mesure
change. Ce TP illustre un empaquetage MSB simple,
pas le format natif d'un CR3 ou d'un ARW.

Pour la compression, comparer l'encodage d'une
plage uniforme et de données pseudo-aléatoires.
Comparer les tailles de fichiers, puis vérifier
les échantillons décodés. Ne pas prendre le ratio
de compression comme un score de qualité visuelle.

\subsection{Observer une perte sans logiciel photographique}
Calculer la DCT d'un bloc constant 12 : le coefficient
DC vaut 96. Quantifier ce coefficient avec $Q=20$
donne 5, puis une reconstruction DC=100. Le bloc
inverse devient 12,5 partout, avant les arrondis
de sortie. L'erreur de 0,5 n'est pas un défaut
de la transformée inverse : elle a été introduite
par la quantification.

\subsection{Protocole comparatif RAW minimal}
Conserver un original RAW, un aperçu extrait si
disponible et deux développements explicitement
paramétrés. Désactiver ou harmoniser l'optique,
le débruitage, la netteté et les réglages ; choisir
le même profil cible et documenter As Shot.
Exporter en TIFF sRGB 16 bits sans changement
de dimensions. Vérifier métadonnées et géométrie
avant de calculer MAE et différences Lab.
Ce protocole rapproche les conditions ; il ne
garantit pas que les moteurs ont le même rendu
de base ni qu'ils sont physiquement fidèles.

\section{Exercices corrigés}
\subsection*{1. Estimer la mémoire}
\emph{Question.} Une photo $8000\times6000$ est chargée
en RGB uint16 puis en RGBA float32. Quelles tailles ?
\emph{Correction.} Il y a 48 millions de pixels.
RGB uint16 : $48\times10^6\times6=288$ MB.
RGBA float32 : $48\times10^6\times16=768$ MB.
Garder les deux représente 1056 MB avant caches
et intermédiaires. La taille compressée n'intervient
pas dans ce calcul.

\subsection*{2. Lire un entier}
\emph{Question.} Que vaut \texttt{00 10} sur 16 bits ?
\emph{Correction.} Big-endian : 16. Little-endian :
4096. Le fichier doit annoncer sa convention.

\subsection*{3. Calculer une adresse}
\emph{Question.} RGB uint16, stride 40, pixel $(3,2)$,
canal 1, origine $A_0$ ?
\emph{Correction.} $A_0+2\times40+3\times3\times2+
1\times2=A_0+100$. Ne pas remplacer le stride
par une largeur supposée.

\subsection*{4. Normaliser un RAW}
\emph{Question.} $D=4096$, $B=1024$, $W=15360$.
\emph{Correction.} $(4096-1024)/(15360-1024)
=3072/14336=3/14\approx0{,}214286$.
Utiliser 65535 au dénominateur parce que le
tampon est uint16 serait une erreur.

\subsection*{5. Appliquer un EV}
\emph{Question.} $L=0{,}18$, réglage $+2$ EV.
\emph{Correction.} $L'=0{,}72$. Encoder seulement
ensuite pour une sortie sRGB ; multiplier
directement le code 0,46136 par quatre écrêterait
abusivement le résultat.

\subsection*{6. Expliquer un hash différent}
\emph{Question.} Deux PNG donnent les mêmes pixels,
mais ont des hashes différents. Est-ce une perte ?
\emph{Correction.} Pas nécessairement. Filtres,
compression, chunks ou métadonnées peuvent différer.
Il faut distinguer identité du fichier et égalité
de ses pixels interprétés.

\subsection*{7. Interpréter le clipping}
\emph{Question.} Le pixel linéaire $(1,0,0)$ est-il noir ?
\emph{Correction.} Non. Il est saturé en rouge dans
la représentation considérée, avec deux canaux
nuls. Un compteur « au moins un canal bas »
ne mesure donc pas le nombre de pixels noirs.

\subsection*{8. Comparer deux TIFF}
\emph{Question.} Même scène, même taille, forte MAE.
Que vérifier avant de corriger les couleurs ?
\emph{Correction.} Orientation, recadrage d'un
pixel, optique, dématriçage, profil ICC, fonction
de transfert, profondeur et normalisation, puis
réglages et rendu de base. Une taille identique
ne démontre pas l'alignement.

\subsection*{9. Comprendre « sans perte »}
\emph{Question.} Un RAW réduit est compressé sans
perte. Peut-on récupérer la résolution originale ?
\emph{Correction.} Pas en général. Le codec peut
conserver exactement une représentation déjà
réduite. La garantie porte sur son entrée.

\subsection*{10. Identifier une fausse conversion}
\emph{Question.} Renommer TIFF en DNG ou assigner
sRGB à des codes Adobe RGB convertit-il les données ?
\emph{Correction.} Non. Le premier change un nom,
le second une interprétation. Une conversion
nécessite des transformations et structures
cohérentes avec le résultat.

\subsection*{11. Diagnostiquer une rotation double}
\emph{Question.} Un export portrait apparaît
tourné encore de 90 degrés. Cause possible ?
\emph{Correction.} Les pixels ont pu être
réorientés sans remettre l'orientation à l'identité.
Vérifier le tableau et la métadonnée séparément.

\subsection*{12. Distinguer précision et dynamique}
\emph{Question.} Convertir un JPEG 8 bits en
TIFF 16 bits récupère-t-il les hautes lumières
perdues ?
\emph{Correction.} Non. Cela augmente la précision
disponible pour les traitements futurs, sans
reconstituer les valeurs déjà écrêtées ou quantifiées.

\section{Glossaire}
\begin{longtable}{p{3.4cm}p{11.2cm}}
\toprule Terme & Définition dans ce cours \\ \midrule
ADC / CAN & Convertisseur analogique-numérique. \\
Aliasing & Repliement de fréquences lié à un échantillonnage insuffisant. \\
Alpha & Information de couverture/opacité ; droit ou prémultiplié. \\
Aperçu & Image déjà développée, souvent embarquée et réduite. \\
Bit / octet & Chiffre binaire / groupe de huit bits. \\
Bit depth & Nombre de bits de représentation par échantillon. \\
Black level & Offset associé au signal noir dans une convention de codes. \\
BMFF & Famille d'organisation de fichiers média en boîtes. \\
Canal & Champ de valeurs : couleur, alpha, profondeur ou autre grandeur. \\
CFA & Mosaïque de filtres colorés du capteur. \\
Chroma & Composantes décrivant une différence colorée dans un modèle donné. \\
Clipping & Écrêtage de valeurs à une limite. \\
Codec & Règles et implémentation de codage/décodage d'une représentation. \\
Conteneur & Structure organisant flux, propriétés et relations. \\
CRC & Contrôle d'intégrité destiné notamment aux erreurs accidentelles. \\
DCP & Profil de caméra de la famille DNG ; matrices et rendu possibles. \\
DCT & Transformée en cosinus discrète. \\
Dématriçage & Estimation de composantes couleur absentes du CFA. \\
DEFLATE & Compression combinant références LZ77 et codage Huffman. \\
Dithering & Signal ajouté pour rendre certaines erreurs de quantification moins structurées. \\
DNG & Format de négatif numérique fondé sur TIFF et enrichi de calibrations. \\
Endianness & Ordre des octets d'une valeur multioctet. \\
Entropie & Quantité moyenne d'information dans un modèle probabiliste. \\
EOTF / OETF & Relations de transfert vers la lumière d'affichage / vers un signal encodé. \\
EV & Unité logarithmique d'exposition ; un EV correspond à un facteur deux. \\
Exif & Conventions de données photographiques et structures associées. \\
FFT & Famille d'algorithmes rapides de transformée de Fourier discrète. \\
Gamut & Ensemble des couleurs représentables ou reproductibles. \\
Gain & Facteur multiplicatif, à distinguer d'une nouvelle acquisition de photons. \\
HDR / SDR & Contextes de dynamique étendue / standard, à préciser selon l'étape. \\
ICC & Famille de profils et contrats de gestion colorimétrique. \\
IFD & Répertoire de champs d'image dans TIFF et familles apparentées. \\
Illuminant & Description spectrale ou normalisée d'un éclairage de référence. \\
IPTC & Conventions éditoriales, descriptives et administratives des médias. \\
JPEG & Famille de procédés de codage ; le mode courant utilise DCT et pertes. \\
LCP & Profil de corrections d'objectif, distinct d'un profil couleur caméra. \\
Linéaire & Proportionnel à une grandeur explicitement définie. \\
Luma / luminance & Signal pondéré souvent non linéaire / grandeur colorimétrique Y. \\
LUT & Table de correspondance, éventuellement multidimensionnelle. \\
MAD & Médiane des écarts absolus à une médiane. \\
MakerNote & Bloc de métadonnées propre au constructeur. \\
Métamérisme & Réponses identiques malgré des spectres différents. \\
MTF / PSF & Transmission des contrastes / étalement spatial d'un point. \\
Offset & Position binaire relative à une origine ; ou décalage additif selon le contexte. \\
Optical-black & Zone de capteur masquée utilisée pour caractériser le noir. \\
PCS & Espace de connexion des profils d'une architecture ICC. \\
Photosite & Élément spatial de mesure du capteur. \\
Pixel & Position/élément d'une représentation matricielle, dont le sens dépend du stade. \\
Planar / interleaved & Composantes en plans séparés / entrelacées. \\
Profil & Description de transformation ; son rôle dépend du type ICC, DCP, etc. \\
Proxy & Représentation réduite pour l'interaction. \\
Quantification & Passage à un ensemble discret de valeurs. \\
RAW & Famille de données photographiques proches des mesures, pas un codec unique. \\
Résidu & Différence entre valeur et prédiction. \\
ROI & Région d'intérêt définie dans un repère et un domaine précis. \\
Round-trip & Encodage puis décodage avec comparaison au domaine initial. \\
Sidecar & Fichier associé portant paramètres ou métadonnées. \\
SNR & Rapport signal/bruit, avec conventions à préciser. \\
sRGB & Espace RGB normalisé avec primaires, blanc et transfert définis. \\
Stride / pitch & Distance mémoire entre deux débuts de ligne. \\
Tag & Champ identifié dans une structure de données. \\
Tone mapping & Adaptation du rendu tonal vers une capacité de sortie. \\
Tuile / bande & Bloc spatial rectangulaire / portion horizontale de stockage. \\
White level & Niveau de référence de saturation ou normalisation, selon le contexte. \\
XMP & Modèle extensible de métadonnées, souvent sérialisé en RDF/XML. \\
XYZ / Lab & Espaces colorimétriques, non simples synonymes de RGB. \\
\bottomrule
\end{longtable}

\section{Bibliographie : quoi lire, dans quel ordre ?}
Pour la structure binaire, commencer par les exemples de
ce cours puis les spécifications PNG, TIFF et Netpbm.
Pour la compression, suivre les calculs de prédiction et
DCT, puis T.81 et les RFC zlib/DEFLATE. Pour la couleur,
lire la fiche sRGB et l'architecture ICC avant les tags
avancés DNG. Pour les capteurs, EMVA fournit un cadre de
mesure plutôt qu'un simple classement commercial.

Les manuels Sony et Nikon de cette bibliographie sont
volontairement attachés à des modèles. Les ressources
Clévy et les décodeurs LibRaw/RawSpeed donnent accès aux
détails des RAW, mais il faut lire leurs limites et leurs
licences. Les branches de code en ligne évoluent : conserver
un commit et les fichiers de test pour une investigation
reproductible de production.

La ressource historique CR2 \url{http://lclevy.free.fr/cr2/}
reste intéressante, mais son accès direct a échoué pendant
la recherche. Le dépôt de l'auteur \texttt{lclevy/libcraw2}
est donc proposé comme point d'entrée alternatif, sans
prétendre avoir relu une page inaccessible. Certains
catalogues de tags ExifTool limitent également l'accès
automatisé ; leur index et les outils locaux permettent
de poursuivre la lecture.

\resultat{Pour comprendre une image, toujours demander : que
mesurent les nombres, comment sont-ils rangés, quelles transformations
ont déjà été faites, et quelles informations faut-il encore appliquer
avant de les afficher ?}
